\section{Attention Variants}\label{sec:exp:attention}\label{sec:exp:results:attention}

Five configurations are trained with the joint objective of \cref{eq:method:loss_total} and otherwise identical hyperparameters (\cref{sec:exp:hparams}).
Vanilla is the joint encoder of \cref{sec:exp:results:vanilla}.
RoPE-TOA applies rotary encoding on the window-local \gls{toa} in the trunk and in both branches.
RoPE-az is the branch-specific rotary design of \cref{sec:method:attention}.
Despite the name, it does not rotate by azimuth or elevation: the trunk and the mode branch rotate by window-local \gls{toa}, and the deinterleaving branch rotates by the two planar Euclidean incidence components $\tilde{u}^{x},\tilde{u}^{y}$ of \cref{eq:method:aoa_scale}.
Bias adds the pairwise physical logits of \cref{eq:method:attn_bias,eq:method:attn_bias_trunk}.
RoPE-TOA+Bias stacks the uniform \gls{toa} rotary layout of RoPE-TOA with that pairwise bias. Incidence is not used as a rotary coordinate.

The rows other than Vanilla test whether explicit physical structure in attention improves clustering relative to standard \gls{mha} (RQ3).
Feature \gls{dbscan} remains the no-encoder control of \cref{tab:exp:vanilla}; it is omitted here.
\Cref{tab:exp:attn_em,tab:exp:attn_md} collect the window-averaged scores.
The best value in each column is in bold; ties are bold on every tied row.

\begin{table}[tbp]
  \centering
  \small
  \caption{Emitter clustering on the test split for jointly trained attention variants. Entries are mean $\pm$ std over the same $960$ windows as \cref{tab:exp:vanilla}. Single training seed.}
  \label{tab:exp:attn_em}
  \begin{tabular}{@{}lccccc@{}}
    \toprule
    Model & $\ARI$ & $\AMI$ & Homogeneity & Completeness & $\lvert\hat{K}-K\rvert$ \\
    \midrule
    Vanilla   & $0.944\pm0.182$ & $0.955\pm0.142$ & $0.989\pm0.054$ & $0.948\pm0.172$ & $0.25\pm0.48$ \\
    RoPE-TOA  & $\mathbf{1.000\pm0.003}$ & $0.998\pm0.007$ & $\mathbf{0.997\pm0.013}$ & $\mathbf{1.000\pm0.000}$ & $0.09\pm0.31$ \\
    RoPE-az   & $0.972\pm0.136$ & $0.976\pm0.113$ & $0.994\pm0.045$ & $0.973\pm0.131$ & $0.14\pm0.37$ \\
    Bias      & $\mathbf{1.000\pm0.003}$ & $\mathbf{0.999\pm0.006}$ & $\mathbf{0.997\pm0.012}$ & $\mathbf{1.000\pm0.000}$ & $\mathbf{0.09\pm0.30}$ \\
    RoPE-TOA+Bias & {---} & {---} & {---} & {---} & {---} \\
    \bottomrule
  \end{tabular}
\end{table}

\begin{table}[tbp]
  \centering
  \small
  \caption{Mode clustering on the test split. Entries follow the same aggregation as \cref{tab:exp:attn_em}.}
  \label{tab:exp:attn_md}
  \begin{tabular}{@{}lccccc@{}}
    \toprule
    Model & $\ARI$ & $\AMI$ & Homogeneity & Completeness & $\lvert\hat{M}-M\rvert$ \\
    \midrule
    Vanilla   & $0.986\pm0.100$ & $0.987\pm0.094$ & $0.993\pm0.037$ & $0.990\pm0.092$ & $0.23\pm0.50$ \\
    RoPE-TOA  & $0.990\pm0.087$ & $0.989\pm0.083$ & $0.995\pm0.023$ & $0.990\pm0.084$ & $0.20\pm0.49$ \\
    RoPE-az   & $0.984\pm0.097$ & $0.987\pm0.084$ & $0.989\pm0.047$ & $0.993\pm0.079$ & $0.23\pm0.50$ \\
    Bias      & $\mathbf{0.993\pm0.074}$ & $\mathbf{0.993\pm0.073}$ & $\mathbf{0.997\pm0.017}$ & $\mathbf{0.994\pm0.073}$ & $\mathbf{0.13\pm0.39}$ \\
    RoPE-TOA+Bias & {---} & {---} & {---} & {---} & {---} \\
    \bottomrule
  \end{tabular}
\end{table}

RoPE-TOA and Bias remove the emitter tail that joint Vanilla leaves in \cref{tab:exp:vanilla,tab:exp:ablation_loss}.
Mean emitter $\ARI$ is $1.000$ to three decimals, $99\%$ of windows sit at or above $0.99$, and none of the $960$ windows falls below $0.90$.
Completeness is $1.000$ with no spread, so a true emitter is not split across clusters.
The leftover errors are rare undercounts, visible as a small homogeneity gap and as the count-error column.
Emitter-only Vanilla already reaches this emitter range in \cref{tab:exp:ablation_loss}, without a usable mode branch.
RoPE-TOA and Bias reach it under the joint objective.

RoPE-az improves on Vanilla for emitters ($\ARI$ $0.972$, $5\%$ of windows below $0.90$) but does not match TOA-only rotary encoding.
Rotating the deinterleaving branch by planar incidence is not, on this test set, a replacement for rotating every layer by \gls{toa}.
Mode scores occupy a narrow band for Vanilla, RoPE-TOA, RoPE-az, and Bias.
Bias is the tightest there, with the smallest count error.

\subsection{Per-window ARI distributions}\label{sec:exp:results:hist}

\Cref{fig:exp:hist_em,fig:exp:hist_md} show the empirical CDF of per-window $\ARI$ for Vanilla, RoPE-TOA, RoPE-az, and Bias.
$\ARI$ is the metric used to select $\varepsilon$ on the validation grid, so the curves correspond to the neural rows of \cref{tab:exp:attn_em,tab:exp:attn_md}.
Each curve is the fraction of the $960$ windows whose $\ARI$ does not exceed the value on the horizontal axis.
The vertical axis is logarithmic so that a tail of a few windows remains visible against the jump at $1$.
$\AMI$, homogeneity, and completeness follow the same right spike and are not repeated as extra pages.

The Vanilla emitter curve leaves the axis well before $0.90$.
RoPE-TOA and Bias stay near zero until just below $1$, and the two curves nearly coincide on this split.
RoPE-az sits between those two shapes.
On the mode branch the four curves rise later and closer together; Bias is the last to leave zero.

\begin{figure}[tbp]
  \centering
  \expfig{hist_em.pdf}
  \caption{Empirical CDF of emitter $\ARI$ over the $960$ non-overlapping test windows. The vertical axis is logarithmic.}
  \label{fig:exp:hist_em}
\end{figure}

\begin{figure}[tbp]
  \centering
  \expfig{hist_md.pdf}
  \caption{Empirical CDF of mode $\ARI$ over the same windows as \cref{fig:exp:hist_em}. The vertical axis is logarithmic.}
  \label{fig:exp:hist_md}
\end{figure}

\subsection{Cluster-count recovery}\label{sec:exp:results:kcorr}

\Cref{fig:exp:kcorr_em,fig:exp:kcorr_md} are heatmaps of predicted versus true cluster counts.
Each cell is the number of test windows with that pair.
The diagonal is exact count recovery.
Mass above the diagonal is extra clusters (splits).
Mass below it is missing clusters (merges).
That is why \cref{tab:exp:attn_em,tab:exp:attn_md} report homogeneity and completeness separately, and why they include a mean absolute count error.

On the emitter branch, Vanilla leaves the diagonal more often than the other three, with splits and merges in roughly equal number.
RoPE-TOA never oversplits an emitter window in this set: the off-diagonal mass is only undercounts.
Bias is the same pattern, up to one window.
RoPE-az is closer to the diagonal than Vanilla and still produces some extra clusters.

The mode heatmaps are diagonal-dominant for every model.
When they miss, they more often undercount than invent labels, and the misses concentrate at the crowded end of the catalogue, where a window can contain a dozen modes (\cref{sec:exp:datasets}).
Bias keeps more mass on the diagonal than Vanilla or RoPE-az.

\begin{figure}[tbp]
  \centering
  \expfig{kcorr_em.pdf}
  \caption{Predicted versus true number of emitters per test window. Cell values are window counts. The line is $\hat{K}=K$.}
  \label{fig:exp:kcorr_em}
\end{figure}

\begin{figure}[tbp]
  \centering
  \expfig{kcorr_md.pdf}
  \caption{Predicted versus true number of modes per test window. Layout as in \cref{fig:exp:kcorr_em}.}
  \label{fig:exp:kcorr_md}
\end{figure}

\subsection{Model size and inference cost}\label{sec:exp:results:efficiency}

\Cref{tab:exp:efficiency} reports trainable parameters and wall-clock time per test window on the A100 of \cref{sec:exp:setup}.
One window is forwarded at a time.
A short GPU warmup is excluded from the mean.
Clustering time is \gls{dbscan} on the embeddings of \cref{eq:method:cluster_decoding}.

Rotary encoding adds no learned weights, so both RoPE models dump the same $16\,081\,536$ parameters as Vanilla.
That is fewer extra parameters than a learned positional table of length $L$, and fewer than the physical bias, which stores one scalar per head, layer, and selected coordinate.
Those scalars do not change the dumped total at the reported precision.
They do change the encoder time: Bias spends $30$~ms in the forward pass where Vanilla spends $1.5$~ms.
The extra work is the pairwise differences $B^{(p)}_{ij}=\lvert p_i-p_j\rvert$ of \cref{eq:method:attn_bias,eq:method:attn_bias_trunk}.
Those matrices depend only on the window coordinates, so they are the same for every head and every attention sub-layer.
The implementation rebuilds them in each head of each sub-layer and then scales by that head's $\lambda$.
A single copy per coordinate, computed once per window, would remove that repetition.
No such cache was used.
The time in \cref{tab:exp:efficiency} includes that redundant construction, not only the add into the logits.
\Gls{dbscan} is about $60$~ms for every neural model, so it dominates Vanilla and both RoPE variants.
Feature \gls{dbscan} has no encoder weights; its clustering time is $21$~ms.
Memory for the pairwise maps is taken up in \cref{sec:disc:interp}.

\begin{table}[tbp]
  \centering
  \caption{Trainable parameters and mean time per test window. Encoder is the forward pass, or the feature copy for Feature \gls{dbscan}. Clustering is \gls{dbscan} on both branches.}
  \label{tab:exp:efficiency}
  \begin{tabular}{@{}lcccc@{}}
    \toprule
    Model & Parameters & Encoder (ms) & Clustering (ms) & Total (ms) \\
    \midrule
    Feature DBSCAN & $0$ & $0.1$ & $20.9$ & $21.0$ \\
    Vanilla   & $16\,081\,536$ & $1.5$ & $61.2$ & $62.7$ \\
    RoPE-TOA  & $16\,081\,536$ & $2.5$ & $59.9$ & $62.4$ \\
    RoPE-az   & $16\,081\,536$ & $2.9$ & $61.3$ & $64.2$ \\
    Bias      & $16\,081\,536$ & $29.6$ & $62.0$ & $91.6$ \\
    RoPE-TOA+Bias & {---} & {---} & {---} & {---} \\
    \bottomrule
  \end{tabular}
\end{table}
